% New introductory planning exercise; no graduate HW06 material copied.
% Shared layout adapted from the instructor's ME 537 homework sources.
\documentclass[12pt,letterpaper]{article}
\usepackage{mathptmx}
\usepackage[T1]{fontenc}
\usepackage{xcolor,amsmath,amssymb,xurl,listings,enumitem,tikz}
\textwidth=6.5in
\textheight=9in
\oddsidemargin=0in
\evensidemargin=0in
\topmargin=0in
\headheight=0in
\headsep=0in
\parskip=0.08in
\parindent=0in
\setlist[enumerate]{itemsep=0.04in,topsep=0.04in}
\lstset{basicstyle=\ttfamily\footnotesize,breaklines=true,columns=fullflexible,
keepspaces=true,showstringspaces=false,aboveskip=0.04in,belowskip=0.04in}
\newcommand{\R}[2]{{}^{#1}\!R_{#2}}
\newcommand{\T}[2]{{}^{#1}\!T_{#2}}
\newcommand{\courseheader}[2]{%
\setlength{\fboxrule}{2pt}\setlength{\fboxsep}{2mm}%
\fcolorbox{black}{white}{\parbox{\dimexpr\linewidth-2\fboxsep-2\fboxrule\relax}{\centering\textbf{Introduction to Robotics Homework \##1 -- Draft}}}\par
\textbf{#2}\par}
\newcommand{\conventions}{Use right-handed rotations, column vectors, radians,
and meters. $\R{A}{B}$ maps coordinates from frame $B$ into frame $A$;
$\T{A}{B}$ also includes the origin of $B$ expressed in $A$.}
\newcommand{\submission}{\textbf{Submit:} labeled hand derivations and sketches,
your changed Python source files, the requested numerical/visual evidence,
and an AI-use statement (including ``none'' if applicable). Show your reasoning
before using code to check hand calculations. Public checks are intentionally
limited. The instructor supplies the deadline and submission destination.}

\begin{document}
\courseheader{6}{Configuration space and RRT-Connect}
\textbf{Schedule:} planning unit; the instructor will set the due date.
Use meters for the base, radians for arm joints, and the provided scenes.
This assignment plans geometric paths; it does not yet assign time to them.

\begin{enumerate}
\item \textbf{Turn a tool goal into a joint goal (IK).}
Tasks are usually given as tool positions, but the planner needs joint
configurations. The provided planar arm (\texttt{arm\_scene()}) has two
revolute joints, lengths $0.4$ and $0.3$ m, and limits $\pm2.8$ rad. Its
tool must reach $p_g=[0.4434,\ 0.3429]$ m.
\begin{enumerate}
\item In \path{src/intro_robotics/kinematics.py}, complete
\texttt{SerialArm.ik\_position}. Repeat: compute the XY
error $e=p_g-p(q)$; stop when $\|e\|$ is below the tolerance; otherwise take
the step $q\leftarrow q+$\texttt{damped\_least\_squares}$(J_{xy}(q),e,\lambda)$
from HW05 and keep each joint within its limits. Return the final $q$ and
whether it converged. Why can the step treat $e$ like a velocity?
\item Run it from $q_0=[0.5,-1]$ and from $q_0=[0,1]$. Report both
solutions, confirm each with your FK, and sketch the two arm shapes. Check
both with \texttt{scene.is\_free(q)}.
\item The start is $q_s=[-1.2,1.3]$. Which solution is closer to it in joint
space? After problem 5, plan to each goal with the same settings and compare
the path lengths. Problems 2--6 use $q_g=[1.2,-1.3]$.
\end{enumerate}

\item \textbf{A path describes the whole robot.}
The same arm faces a circular obstacle centered at $(0.55,0)$ with radius
$0.07$ m. Each link has radius $0.015$ m. Its start and goal configurations
are $q_s=[-1.2,1.3]$ and $q_g=[1.2,-1.3]$ radians.
\begin{enumerate}
\item Sketch the start and goal arm poses and calculate their tool positions
using your FK code. Identify the two configuration-space coordinates.
\item The midpoint of the straight joint-space interpolation is $[0,0]$.
Show by hand why the second link collides there even though the tool point
itself is outside the obstacle. Check the start and goal using the supplied
\texttt{arm\_scene()} and \texttt{scene.is\_free(q)}.
\item Explain why valid start and goal configurations do not by themselves
prove that the motion between them is collision-free.
\end{enumerate}

\end{enumerate}

\newpage
\begin{enumerate}[start=3]
\item \textbf{Trace two trees by hand.}
In an empty 2D configuration space, the start tree contains a root at
$(0,0)$, node 1 at $(0.2,0)$ with parent 0, and node 2 at $(0.2,0.2)$ with
parent 1. The goal tree has only its root at $(0.8,0.2)$. The planner tries
to connect the goal tree to node 2 with step size $0.2$.
\begin{enumerate}
\item List the new goal-tree nodes and their parent indices. Assume exact
arithmetic and that all edges are valid.
\item List the complete path in start-to-goal order. Include the meeting
point once. Explain which part of the goal-tree path had to be reversed.
\end{enumerate}

\item \textbf{Implement one tree extension.}\par
In \path{src/intro_robotics/rrt.py}, complete \texttt{extend\_tree}.
The supplied \texttt{Tree} stores configurations in \texttt{points} and
matching parent indices in \texttt{parents}; \texttt{path\_to} follows
parents back to the root. Input checks and \texttt{edge\_is\_free} are supplied.
\begin{enumerate}
\item Find the nearest existing node and move toward the target by at most
\texttt{step}. Check the entire connecting edge at the specified resolution.
\item For a valid extension, append the new point and its parent; return
its index and whether it reached the target. For a blocked edge return
\texttt{(None, False)} without changing the tree. If the target is already a
node, return that node and \texttt{True}.
\end{enumerate}
\end{enumerate}

\newpage
\begin{enumerate}[start=5]
\item \textbf{Connect the trees.}
Complete \texttt{rrt\_connect} in the same file. Endpoint checks, a seeded
random generator, and the tree class are provided.
\begin{enumerate}
\item Start one tree at each endpoint. Sample inside the supplied bounds,
extend the active tree once, and repeatedly extend the other tree toward the
new point until it connects or is blocked. Alternate active trees each attempt.
For example, the provided generator can sample with:
\begin{lstlisting}
sample = rng.uniform(bounds[:, 0], bounds[:, 1])
\end{lstlisting}
\item When connected, use \texttt{path\_to} for each tree and assemble a path
from start to goal without duplicating the meeting point. Return \texttt{None}
if the iteration budget is exhausted. Explain why a failure does not prove
that no path exists.
\end{enumerate}

\item \textbf{Check and visualize your planner.}
From \path{/workspace} in the course container, run:
\begin{lstlisting}
python -m homework.hw06.hw06_planning
python -m homework.hw06.hw06_planning --show
python -m homework.hw06.hw06_planning --save tmp/hw06.png
\end{lstlisting}
The first command runs limited public checks. The viewer shows a circular
base, the arm's configuration-space path, and several physical arm poses.
\begin{enumerate}
\item In a short script, call \texttt{base\_scene()} and run the planner for
at least three seeds using the same step, edge resolution, and iteration
budget. For each run report success/failure, planning time, and the sum of
Euclidean distances between adjacent path points when a path is found.
\item Check every path edge with the supplied function at a finer resolution.
Compare the routes and report any failure. Why is this still not a proof of
continuous collision freedom?
\item Save the supplied arm view. Explain how a point in its configuration
plot corresponds to an entire physical arm pose. State one limitation of
the obstacle, joint metric, or collision model.
\end{enumerate}
\end{enumerate}

\textbf{What to understand:} sampling and connecting two trees search for a
feasible geometric route. The arm's full geometry and connecting motions
must be checked. Time, velocity limits, and path optimization are later work.

\submission
\end{document}
